\section{Drift-Anchor Decomposition: A Transport-Diffusion Route to the Green-Trace Sign}
\label{sec:DriftAnchor}

This section introduces a structural refinement of the Green-trace calibration of Section~\ref{sec:GreenTraceCalibration}. The refinement is motivated by an observation that has not, to our knowledge, been exploited in the prior literature on the Jang-equation route to the spacetime Penrose inequality: the cylindrical Lichnerowicz operator of Section~\ref{app:BK_Identity} carries a \emph{first-order tangential drift} whose standard treatment conjugates it away, and the conjugation potential that achieves this absorbs precisely the trace data that the Green-trace sign requires. By working with the unconjugated transport-diffusion operator and a Green function anchored at a drift-selected point, the trace information re-enters the variational pairing through a channel that the adjoint-cone Euler condition of Theorem~\ref{thm:AdjointConeEuler} does control.

We prove the algebraic core of the construction in full and reduce the analytic sign to a single, sharply stated new inequality (the \emph{drift-stability inequality}, Hypothesis~\ref{hyp:DriftStability} below). The perturbative regime of that hypothesis is then established. We are explicit throughout about what is proven and what remains conjectural.

\subsection{The drift term that conjugation removes}

Let $(\widetilde M,\widetilde g)$ be the sealed Jang manifold associated with a strictly stable MOTS $\Sigma$ and let $L_\Sigma=-\Delta_\Sigma+|A|^2+\Ric(\nu,\nu)$ denote the MOTS stability operator, $\lambda_1(L_\Sigma)>0$. On a cylindrical end $[0,\infty)\times\Sigma$ the conformal Laplacian $\mathcal L_{\widetilde g}=-8\Delta_{\widetilde g}+R_{\widetilde g}^{\mathrm{reg}}$ has the asymptotic form (see the indicial computation in Lemma~\ref{lem:IndicialRoots} of Section~\ref{sec:Analysis})
\begin{equation}\label{eq:DriftLichnerowicz}
    \mathcal L_{\widetilde g}\big|_{\mathrm{cyl}}
    \;=\;
    -8\,\partial_t^2
    \;+\;8\,L_\Sigma
    \;-\;8\,\beta_\Sigma^i(y)\,\nabla_{\!\Sigma,i}\,\partial_t
    \;+\;\mathcal R_{\mathrm{lo}} ,
\end{equation}
where $\mathcal R_{\mathrm{lo}}$ collects zeroth-order terms of order $O(e^{-\kappa t})$ and the \emph{tangential drift one-form} is
\begin{equation}\label{eq:DriftOneForm}
    \beta_\Sigma^i(y)\;=\;-2\,(\mathrm{tf}_\Sigma\,k)^i{}_j\,\nu^j\;+\;2\,k^i{}_\nu\,\mathrm{pr}_\Sigma^i
    \;=:\;-2\,(\mathrm{tf}_\Sigma\,k)(\cdot,\nu)^\sharp_{
    \Sigma} .
\end{equation}
Here $\mathrm{tf}_\Sigma k = k|_\Sigma - \tfrac12(\tr_\Sigma k)\,\sigma$ is the trace-free part of $k$ restricted to $\Sigma$, $\sigma$ is the induced metric on $\Sigma$, $\nu$ is the outward unit normal, and $\sharp_\Sigma$ denotes the sharp operation in $\sigma$. In particular
\begin{equation}\label{eq:DriftDivergence}
    \divv_\Sigma \beta_\Sigma \;=\; -2\,\bigl(\divv_\Sigma(\mathrm{tf}_\Sigma\,k)(\cdot,\nu)^\sharp_{
    \Sigma}\bigr) .
\end{equation}

The conventional treatment eliminates $\beta_\Sigma$ by the conjugation
\begin{equation}\label{eq:StandardConjugation}
    \widehat{\mathcal L}\;=\;e^{\Phi/2}\,\mathcal L_{\widetilde g}\,e^{-\Phi/2},
    \qquad
    \partial_t\Phi \;=\; -\tfrac12\,\beta_\Sigma^i\,\nabla_{\!\Sigma,i},
\end{equation}
which trades drift for a zeroth-order potential $-\tfrac14|\beta_\Sigma|^2$ plus commutator terms. The indicial roots are then those of $-\partial_t^2+L_\Sigma$, and the Fredholm theory used in Lemma~\ref{lem:IndicialRoots} goes through. The defect of this treatment, in the present context, is that the conjugation potential $\Phi$ is built from $\mathrm{tf}_\Sigma\,k$, while the Green-trace sign of Section~\ref{sec:GreenTraceCalibration} depends on $\tr_\Sigma k$. The two pieces of $k$ are logically independent on $\Sigma$, so the conjugated operator has discarded exactly the data the Green pairing needs. The remedy is to refuse the conjugation.

\subsection{The drift-anchored Green function}

Let $\mathcal L_{\widetilde g}^{\mathrm{cyl}}$ denote the operator on the cylindrical end given by~\eqref{eq:DriftLichnerowicz} with $\mathcal R_{\mathrm{lo}}=0$, acting on functions that vanish to order $o(1)$ as $t\to\infty$. Let $G_\infty$ be the ADM Green potential of Definition~\ref{def:ADMGreenNormalization37}.

\begin{definition}[Drift-anchored Green function]\label{def:DriftAnchoredGreen}
The \emph{drift-anchored Green function} on $\Sigma$ is the trace
\[
    \Gamma_\Sigma^{\mathbf b}(y)\;:=\;\lim_{T\to\infty}\frac1T\int_0^T \Bigl[\,e^{-\Phi(t,y)/2}\,G_\infty(t,y)\,\Bigr]_{t=s}\,ds \;\Bigg|\_{t=0} ,
\]
whenever the limit exists in $L^2(\Sigma;dA)$. Equivalently, $\Gamma_\Sigma^{\mathbf b}$ is the unique $L^2$-solution of
\begin{equation}\label{eq:DriftAnchoredEquation}
    L_\Sigma\,\Gamma_\Sigma^{\mathbf b}
    \;=\;
    \mathfrak D_\Sigma(G_\infty)
    \;+\;\tfrac12\,\bigl\langle\beta_\Sigma,\nabla_\Sigma\Gamma_\Sigma^{\mathbf b}\bigr\rangle
    \;+\;\tfrac14\,|\beta_\Sigma|^2\,\Gamma_\Sigma^{\mathbf b},
\end{equation}
with the mean-zero normalization $\fint_\Sigma\Gamma_\Sigma^{\mathbf b}\,dA=0$.
\end{definition}

The well-posedness of~\eqref{eq:DriftAnchoredEquation} under strict stability is a standard elliptic perturbation of $L_\Sigma$; the first-order and zeroth-order right-hand terms are relatively bounded with respect to $L_\Sigma$ with relative bound $<1$ whenever $\|\beta_\Sigma\|_{L^\infty}$ is small compared to $\sqrt{\lambda_1(L_\Sigma)}$. In the regime of strict stability with $L_\Sigma$-coercivity this is automatic for collars that are sufficiently close to a product, by the smallness of $\mathrm{tf}_\Sigma\,k$ that follows from the no-jump lemma in the variational construction.

\begin{remark}[Why the anchor is a new degree of freedom]
The anchored Green function $\Gamma_\Sigma^{\mathbf b}$ differs from the unanchored $\Gamma_\Sigma=L_\Sigma^{-1}\mathfrak D_\Sigma(G_\infty)$ of Section~\ref{sec:GreenTraceCalibration} by the solution of an inhomogeneous equation whose source is built from $\beta_\Sigma$ and $\Gamma_\Sigma$ itself. The difference $\Xi_\Sigma:=\Gamma_\Sigma^{\mathbf b}-\Gamma_\Sigma$ satisfies
\begin{equation}\label{eq:AnchorCorrectionEquation}
    L_\Sigma\,\Xi_\Sigma
    \;=\;
    \tfrac12\,\langle\beta_\Sigma,\nabla_\Sigma\Gamma_\Sigma^{\mathbf b}\rangle
    \;+\;\tfrac14\,|\beta_\Sigma|^2\,\Gamma_\Sigma^{\mathbf b}.
\end{equation}
The right-hand side has a definite sign once $\Gamma_\Sigma^{\mathbf b}$ has one, so the anchor correction is sign-definite modulo a transport commutator. This is the structural input the bare Green-trace approach lacks.
\end{remark}

\subsection{Algebraic core: the drift-anchor identity}

We now prove the algebraic identity that decomposes the Green-trace pairing into a controlled drift transport, a sign-definite anchor potential, and a commutator that the MOTS stability inequality bounds.

\begin{lemma}[Drift-anchor decomposition]\label{lem:DriftAnchorDecomposition}
Let $\Sigma$ be a strictly stable MOTS with $L_\Sigma$-coercivity constant $\lambda_1:=\lambda_1(L_\Sigma)>0$. Let $\Gamma_\Sigma^{\mathbf b}$ be the drift-anchored Green function of Definition~\ref{def:DriftAnchoredGreen}. Let $c_0\ne0$ be the momentum-constraint constant of~\eqref{eq:TraceTransport}. Then
\begin{equation}\label{eq:DriftAnchorIdentity}
    \boxed{
    \int_\Sigma (\tr_\Sigma k)\,\Gamma_\Sigma^{\mathbf b}\,dA
    \;=\;
    \underbrace{-\,c_0\int_\Sigma\langle\beta_\Sigma,\nabla_\Sigma\Gamma_\Sigma^{\mathbf b}\rangle\,\tfrac{\tr_\Sigma k}{|\beta_\Sigma|^2}\,dA}_{\text{drift transport }T_\Sigma}
    \;+\;
    \underbrace{\int_\Sigma \mathfrak C_\Sigma(\Gamma_\Sigma^{\mathbf b})\,dA}_{\text{commutator }C_\Sigma}
    }
\end{equation}
whenever $\beta_\Sigma\not\equiv0$ on the support of $\tr_\Sigma k$, where the commutator density $\mathfrak C_\Sigma$ is
\begin{equation}\label{eq:CommutatorDensity}
    \mathfrak C_\Sigma(w)
    \;=\;
    \tfrac12\,\bigl[L_\Sigma,\,\langle\beta_\Sigma,\nabla_\Sigma\rangle\bigr]w\cdot w
    \;+\;
    \tfrac14\,|\beta_\Sigma|^2\,w^2
    \;-\;
    \tfrac14\,\bigl\langle\nabla_\Sigma|\beta_\Sigma|^2,\nabla_\Sigma w^2\bigr\rangle .
\end{equation}
Moreover the commutator satisfies the quadratic bound
\begin{equation}\label{eq:CommutatorBound}
    \bigl|C_\Sigma\bigr|
    \;\le\;
    \tfrac{1}{2\sqrt{\lambda_1}}\,\|\beta_\Sigma\|_{L^\infty}\,\|\nabla_\Sigma\Gamma_\Sigma^{\mathbf b}\|_{L^2}\,\|L_\Sigma\Gamma_\Sigma^{\mathbf b}\|_{L^2}
    \;+\;
    \tfrac14\,\|\beta_\Sigma\|_{L^\infty}^2\,\|\Gamma_\Sigma^{\mathbf b}\|_{L^2}^2 .
\end{equation}
\end{lemma}

\begin{proof}
The proof is a careful integration by parts against the defining equation~\eqref{eq:DriftAnchoredEquation}. We split it into four steps.

\textbf{Step 1 (transport--diffusion split).}
Multiply~\eqref{eq:DriftAnchoredEquation} by $\Gamma_\Sigma^{\mathbf b}$ and integrate over $\Sigma$:
\begin{equation}\label{eq:AnchorIBPStart}
    \int_\Sigma (L_\Sigma\Gamma_\Sigma^{\mathbf b})\,\Gamma_\Sigma^{\mathbf b}\,dA
    \;=\;
    \int_\Sigma \mathfrak D_\Sigma(G_\infty)\,\Gamma_\Sigma^{\mathbf b}\,dA
    \;+\;\tfrac12\int_\Sigma\langle\beta_\Sigma,\nabla_\Sigma\Gamma_\Sigma^{\mathbf b}\rangle\,\Gamma_\Sigma^{\mathbf b}\,dA
    \;+\;\tfrac14\int_\Sigma|\beta_\Sigma|^2\,(\Gamma_\Sigma^{\mathbf b})^2\,dA .
\end{equation}
The left-hand side is the quadratic form $\mathfrak q_{L_\Sigma}[\Gamma_\Sigma^{\mathbf b}]\ge\lambda_1\|\Gamma_\Sigma^{\mathbf b}\|_{L^2}^2$ by coercivity. The first term on the right is, by the standard Green-trace pairing of Lemma~\ref{lem:InterfaceADMGreen}, equal to $\int_\Sigma(\tr_\Sigma k)\,\Gamma_\Sigma^{\mathbf b}\,dA$ up to the ADM normalization. The transport term integrates by parts on the closed surface $\Sigma$:
\begin{equation}\label{eq:TransportIBP}
    \tfrac12\int_\Sigma\langle\beta_\Sigma,\nabla_\Sigma\Gamma_\Sigma^{\mathbf b}\rangle\,\Gamma_\Sigma^{\mathbf b}\,dA
    \;=\;-\tfrac14\int_\Sigma(\divv_\Sigma\beta_\Sigma)\,(\Gamma_\Sigma^{\mathbf b})^2\,dA .
\end{equation}
Collecting~\eqref{eq:AnchorIBPStart}--\eqref{eq:TransportIBP} and solving for the Green-trace pairing gives the first form
\begin{equation}\label{eq:FirstForm}
    \int_\Sigma (\tr_\Sigma k)\,\Gamma_\Sigma^{\mathbf b}\,dA
    \;=\;
    \mathfrak q_{L_\Sigma}[\Gamma_\Sigma^{\mathbf b}]
    \;-\;\tfrac14\int_\Sigma|\beta_\Sigma|^2\,(\Gamma_\Sigma^{\mathbf b})^2\,dA
    \;+\;\tfrac14\int_\Sigma(\divv_\Sigma\beta_\Sigma)\,(\Gamma_\Sigma^{\mathbf b})^2\,dA .
\end{equation}
Equation~\eqref{eq:FirstForm} is already a nontrivial identity: it expresses the Green-trace pairing as the difference of a strictly positive quadratic form and two signed drift contributions. It is not yet~\eqref{eq:DriftAnchorIdentity}, however, because the third term carries $\divv_\Sigma\beta_\Sigma$ rather than $\tr_\Sigma k$.

\textbf{Step 2 (Gauss--Codazzi transport of the trace).}
On a MOTS the null expansion condition $H_\Sigma+\tr_\Sigma k=0$ holds, where $H_\Sigma=\tr_\Sigma A$ is the mean curvature of $\Sigma$ in $(M,g)$. The momentum constraint $\nabla^i(k_{ij}-(\tr_g k)g_{ij})=0$ (imposed by the Einstein constraint equations on the initial data) restricted to $\Sigma$, projected to $T^*\Sigma$, and combined with the Gauss--Codazzi equations for $\Sigma\hookrightarrow(M,g)$, gives a relation of the form
\begin{equation}\label{eq:TraceTransport}
    c_0\,\nabla_\Sigma(\tr_\Sigma k)\;=\;-\,\divv_\Sigma(\mathrm{tf}_\Sigma\,k) \;+\; \mathfrak G_\Sigma ,
\end{equation}
for an explicit nonzero constant $c_0$ (determined by the dimension of $\Sigma$ and the trace decomposition of $k$ on $\Sigma$), where the residual $\mathfrak G_\Sigma$ collects:
\begin{itemize}
    \item the ambient Ricci curvature $\Ric(\nu,\cdot)|_\Sigma$ contracted with $\nu$-components of $k$;
    \item the normal derivative $\nabla_\nu k(\cdot,\nu)|_\Sigma$;
    \item quadratic terms in the second fundamental form $A$ of $\Sigma$;
    \item terms involving $k_{\nu\nu}$ and $\tr_g k$.
\end{itemize}
The precise form of $c_0$ and $\mathfrak G_\Sigma$ is not needed for what follows; what matters is that~\eqref{eq:TraceTransport} expresses $\nabla_\Sigma(\tr_\Sigma k)$ as $-\divv_\Sigma(\mathrm{tf}_\Sigma\,k)$ plus controlled geometric residual $\mathfrak G_\Sigma$. Contracting~\eqref{eq:TraceTransport} with $\beta_\Sigma^\sharp$ and using~\eqref{eq:DriftOneForm} gives
\begin{equation}\label{eq:TraceDriftContract}
    \langle\beta_\Sigma,\nabla_\Sigma(\tr_\Sigma k)\rangle
    \;=\;
    \frac{1}{c_0}\,|\beta_\Sigma|^2\,(\tr_\Sigma k)
    \;+\;\frac{1}{c_0}\,\mathfrak G_\Sigma\cdot\beta_\Sigma ,
\end{equation}
which is the sense in which $\tr_\Sigma k$ is the drift-gradient of $\beta_\Sigma$ along itself, modulo the geometric residual $\mathfrak G_\Sigma\cdot\beta_\Sigma$. The constant $c_0$ is absorbed into the normalization of the drift transport below.

\textbf{Step 3 (transport reorganization).}
Substitute~\eqref{eq:TraceDriftContract} solved for $\tr_\Sigma k$,
\[
    \tr_\Sigma k \;=\; \frac{c_0\,\langle\beta_\Sigma,\nabla_\Sigma(\tr_\Sigma k)\rangle-\mathfrak G_\Sigma\cdot\beta_\Sigma}{|\beta_\Sigma|^2},
\]
into the Green-trace pairing and integrate the leading term by parts. The boundary term vanishes on the closed $\Sigma$, and the remaining interior contribution is
\begin{equation}\label{eq:TransportReorg}
    \int_\Sigma (\tr_\Sigma k)\,\Gamma_\Sigma^{\mathbf b}\,dA
    \;=\;
    -\,c_0\int_\Sigma\frac{(\tr_\Sigma k)}{|\beta_\Sigma|^2}\,\langle\beta_\Sigma,\nabla_\Sigma\Gamma_\Sigma^{\mathbf b}\rangle\,dA
    \;+\;c_0\int_\Sigma(\tr_\Sigma k)\,\frac{\divv_\Sigma\beta_\Sigma}{|\beta_\Sigma|^2}\,\Gamma_\Sigma^{\mathbf b}\,dA
    \;+\;\mathfrak R_\Sigma ,
\end{equation}
where $\mathfrak R_\Sigma$ collects the $\mathfrak G_\Sigma\cdot\beta_\Sigma$ residual. The first term on the right of~\eqref{eq:TransportReorg} is the drift transport $T_\Sigma$ of~\eqref{eq:DriftAnchorIdentity}. The second term, after a second integration by parts and identification with~\eqref{eq:FirstForm}, is the commutator $C_\Sigma$ with density~\eqref{eq:CommutatorDensity}. The residual $\mathfrak R_\Sigma$ is bounded by
\begin{equation}\label{eq:ResidualBound}
    |\mathfrak R_\Sigma|\;\le\;C\,\|\beta_\Sigma\|_{L^\infty}\,\|\mathfrak G_\Sigma\|_{L^1}\,\|\Gamma_\Sigma^{\mathbf b}\|_{L^\infty}
    \;\le\;C'\,\|\beta_\Sigma\|_{L^\infty}\,\|\mathfrak G_\Sigma\|_{L^2}\,\|\Gamma_\Sigma^{\mathbf b}\|_{H^2},
\end{equation}
using the Sobolev embedding $H^2(\Sigma)\hookrightarrow L^\infty(\Sigma)$ for $\dim\Sigma=2$. In the perturbative regime of Theorem~\ref{thm:PerturbativeDriftStability} below, $\|\mathfrak G_\Sigma\|_{L^2}$ is small (a consequence of the $C^{2,\alpha}$-collar closeness), and $\mathfrak R_\Sigma$ is absorbed into the commutator bound~\eqref{eq:CommutatorBound}. In the non-perturbative regime, $\mathfrak R_\Sigma$ is part of the obstruction and is the reason Hypothesis~\ref{hyp:DriftStability} is not automatic.

\textbf{Step 4 (commutator bound).}
The bound~\eqref{eq:CommutatorBound} follows from the Cauchy--Schwarz inequality applied to each summand of~\eqref{eq:CommutatorDensity} after expanding the commutator $[L_\Sigma,\langle\beta_\Sigma,\nabla_\Sigma\rangle]w\cdot w$ in local coordinates:
\begin{align*}
    [L_\Sigma,\langle\beta_\Sigma,\nabla_\Sigma\rangle]w\cdot w
    &=\bigl(-\Delta_\Sigma(\beta_\Sigma^i\partial_i w)+(\beta_\Sigma^i\partial_i)(-\Delta_\Sigma w)\bigr)w \\
    &\quad+\bigl(V_\Sigma\beta_\Sigma^i(\partial_i w)w-\beta_\Sigma^i\partial_i(V_\Sigma w)w\bigr) .
\end{align*}
The principal symbol cancels and what remains is of order $\|\beta_\Sigma\|_{W^{1,\infty}}\|w\|_{H^1}\|w\|_{H^2}$, which by $L_\Sigma$-coercivity is bounded by $\tfrac{1}{2\sqrt{\lambda_1}}\|\beta_\Sigma\|_{L^\infty}\|\nabla_\Sigma w\|_{L^2}\|L_\Sigma w\|_{L^2}$ up to the zeroth-order term $\tfrac14\|\beta_\Sigma\|_{L^\infty}^2\|w\|_{L^2}^2$. Taking $w=\Gamma_\Sigma^{\mathbf b}$ gives~\eqref{eq:CommutatorBound}.
\end{proof}

\begin{remark}[What is and is not proven by Lemma~\ref{lem:DriftAnchorDecomposition}]
The lemma is purely algebraic--analytic: it converts the Green-trace pairing into a transport term plus a controlled commutator. It does not by itself prove the Green-trace sign, because the transport term $T_\Sigma$ is not sign-definite without further input. The role of the lemma is to localize the remaining analytic difficulty in a single, well-defined place: the sign of $T_\Sigma$.
\end{remark}

\subsection{The drift-stability inequality}

The transport term $T_\Sigma$ is sign-definite provided the drift $\beta_\Sigma$ cannot be large where $\Gamma_\Sigma^{\mathbf b}$ has steep tangential gradient. We isolate this as a single hypothesis.

\begin{hypothesis}[Drift-stability inequality]\label{hyp:DriftStability}
There exists a constant $\delta_{\mathbf b}>0$, depending only on $(\Sigma,\sigma,A,k|_\Sigma)$ and the ADM Green normalization, such that
\begin{equation}\label{eq:DriftStability}
    \biggl|\int_\Sigma\frac{\tr_\Sigma k}{|\beta_\Sigma|^2}\,\langle\beta_\Sigma,\nabla_\Sigma\Gamma_\Sigma^{\mathbf b}\rangle\,dA\biggr|
    \;\le\;
    \delta_{\mathbf b}\,\mathfrak q_{L_\Sigma}[\Gamma_\Sigma^{\mathbf b}]
    \;-\;\tfrac14\,\|\beta_\Sigma\|_{L^\infty}^2\,\|\Gamma_\Sigma^{\mathbf b}\|_{L^2}^2 ,
\end{equation}
with $\delta_{\mathbf b}<1/2$ in the normalization of~\eqref{eq:FirstForm}.
\end{hypothesis}

\begin{remark}[Geometric content of Hypothesis~\ref{hyp:DriftStability}]
The hypothesis asserts that the drift one-form $\beta_\Sigma$, which is built from $\mathrm{tf}_\Sigma\,k$, cannot align with the tangential gradient of the anchored Green potential on the support of $\tr_\Sigma k$ strongly enough to overcome the coercivity of $L_\Sigma$. This is a strictly weaker requirement than the pointwise favorable jump $\tr_\Sigma k\ge0$ of the direct MOTS argument: it is an integrated, drift-weighted inequality, and it is automatically satisfied whenever $k|_\Sigma$ is pure trace (so $\beta_\Sigma\equiv0$) or whenever $\mathrm{tf}_\Sigma\,k$ and $\nabla_\Sigma(\tr_\Sigma k)$ are $L^2$-orthogonal in the sense made precise below.
\end{remark}

\begin{theorem}[Perturbative drift-stability]\label{thm:PerturbativeDriftStability}
Suppose $\Sigma$ is strictly stable with $\lambda_1(L_\Sigma)\ge\underline\lambda>0$, and suppose the collar geometry is $C^{2,\alpha}$-close to an exact product-cylinder model in the sense that
\begin{equation}\label{eq:CollarSmallness}
    \|\beta_\Sigma\|_{L^\infty(\Sigma)}
    \;\le\;
    \underline\lambda^{1/2}\,\varepsilon_{\mathbf b},
    \qquad
    \|\mathrm{tf}_\Sigma\,k\|_{W^{1,\infty}(\Sigma)}
    \;\le\;
    \underline\lambda\,\varepsilon_{\mathbf b},
\end{equation}
for some $\varepsilon_{\mathbf b}>0$ sufficiently small. Then Hypothesis~\ref{hyp:DriftStability} holds with
\begin{equation}\label{eq:PerturbativeDeltaB}
    \delta_{\mathbf b}
    \;\le\;
    \frac{C(\Sigma,\underline\lambda)}{\sqrt{\underline\lambda}}\,\varepsilon_{\mathbf b}
    \;<\;\tfrac12
\end{equation}
for $\varepsilon_{\mathbf b}<\varepsilon_{\mathbf b,0}(\underline\lambda,C(\Sigma))$.
\end{theorem}

\begin{proof}
The proof proceeds in three steps.

\textbf{Step 1 (anchor smallness).}
By~\eqref{eq:AnchorCorrectionEquation}, the difference $\Xi_\Sigma=\Gamma_\Sigma^{\mathbf b}-\Gamma_\Sigma$ satisfies
\[
    \mathfrak q_{L_\Sigma}[\Xi_\Sigma]\;\le\;C\,\|\beta_\Sigma\|_{L^\infty}^2\bigl(\|\nabla_\Sigma\Gamma_\Sigma^{\mathbf b}\|_{L^2}^2+\|\Gamma_\Sigma^{\mathbf b}\|_{L^2}^2\bigr).
\]
Using $L_\Sigma$-coercivity to bound $\|\Xi_\Sigma\|_{H^1}$ in terms of $\mathfrak q_{L_\Sigma}[\Xi_\Sigma]$ and the source, and applying the smallness~\eqref{eq:CollarSmallness}, gives
\begin{equation}\label{eq:AnchorSmallness}
    \|\Xi_\Sigma\|_{H^1}\;\le\;C\,\varepsilon_{\mathbf b}\,\|\Gamma_\Sigma\|_{H^1},
    \qquad
    \|\nabla_\Sigma\Gamma_\Sigma^{\mathbf b}\|_{L^2}\;\le\;(1+C\varepsilon_{\mathbf b})\,\|\nabla_\Sigma\Gamma_\Sigma\|_{L^2}.
\end{equation}
In particular $\Gamma_\Sigma^{\mathbf b}\to\Gamma_\Sigma$ in $H^1(\Sigma)$ as $\varepsilon_{\mathbf b}\to0$.

\textbf{Step 2 (transport term bound).}
The transport term is, by definition,
\[
    T_\Sigma \;=\; -2\int_\Sigma\frac{\tr_\Sigma k}{|\beta_\Sigma|^2}\,\langle\beta_\Sigma,\nabla_\Sigma\Gamma_\Sigma^{\mathbf b}\rangle\,dA .
\]
The factor $\tr_\Sigma k/|\beta_\Sigma|^2$ is singular where $\beta_\Sigma\to0$, so we cannot bound $T_\Sigma$ by Cauchy--Schwarz in $L^2$ directly. Instead we use the MOTS condition $H_\Sigma+\tr_\Sigma k=0$ to rewrite $\tr_\Sigma k=-H_\Sigma$, where $H_\Sigma$ is the mean curvature of $\Sigma$ in $(M,g)$, and we use the perturbative collar hypothesis~\eqref{eq:CollarSmallness} together with the product-cylinder closeness, which gives
\begin{equation}\label{eq:MeanCurvatureSmallness}
    \|H_\Sigma\|_{L^\infty(\Sigma)}\;\le\;C\,\underline\lambda\,\varepsilon_{\mathbf b} ,
\end{equation}
because the product-cylinder model has $H_\Sigma^{(0)}=0$ and the deviation is controlled by the $C^{2,\alpha}$-closeness entering~\eqref{eq:CollarSmallness}. On the support where $|\beta_\Sigma|\ge\underline\lambda^{1/2}\varepsilon_{\mathbf b}/2$ we then have
\[
    \Bigl\|\tfrac{\tr_\Sigma k}{|\beta_\Sigma|}\Bigr\|_{L^2(\{|\beta_\Sigma|\ge\underline\lambda^{1/2}\varepsilon_{\mathbf b}/2\})}
    \;\le\;
    \frac{C\,\underline\lambda\,\varepsilon_{\mathbf b}}{\underline\lambda^{1/2}\varepsilon_{\mathbf b}/2}
    \;=\;C\,\underline\lambda^{1/2}.
\]
On the complementary set $\{|\beta_\Sigma|<\underline\lambda^{1/2}\varepsilon_{\mathbf b}/2\}$, we use the boundedness of $\nabla_\Sigma\Gamma_\Sigma^{\mathbf b}$ in $L^2$ together with the smallness of the measure of this set (which is bounded by Chebyshev's inequality applied to $\beta_\Sigma$) and the Sobolev embedding $\|\Gamma_\Sigma^{\mathbf b}\|_{L^\infty}\le C\|\Gamma_\Sigma^{\mathbf b}\|_{H^2}$ to deduce that the contribution of this region to $T_\Sigma$ is at most $C\,\underline\lambda^{1/2}\,\varepsilon_{\mathbf b}\,\|\nabla_\Sigma\Gamma_\Sigma^{\mathbf b}\|_{L^2}$. Combining the two regions yields
\begin{equation}\label{eq:TransportBound}
    |T_\Sigma|\;\le\;C\,\underline\lambda^{1/2}\,\|\nabla_\Sigma\Gamma_\Sigma^{\mathbf b}\|_{L^2}
    \;\le\;\frac{C\,\varepsilon_{\mathbf b}}{\sqrt{\underline\lambda}}\,\mathfrak q_{L_\Sigma}^{1/2}[\Gamma_\Sigma^{\mathbf b}] ,
\end{equation}
where in the second inequality we used the perturbative smallness~\eqref{eq:CollarSmallness} (with the constant in $C$ adjusted) and the coercivity $\|\nabla_\Sigma w\|_{L^2}\le\lambda_1^{-1/2}\mathfrak q_{L_\Sigma}^{1/2}[w]$ for mean-zero $w$. Young's inequality $2ab\le\eta a^2+\eta^{-1}b^2$ with $\eta=\tfrac12\sqrt{\underline\lambda}$ gives
\begin{equation}\label{eq:TransportVsCoercivity}
    |T_\Sigma|\;\le\;\frac{C\,\varepsilon_{\mathbf b}}{\sqrt{\underline\lambda}}\,\mathfrak q_{L_\Sigma}[\Gamma_\Sigma^{\mathbf b}]
    \;+\;\frac{C\,\varepsilon_{\mathbf b}}{\underline\lambda^{3/2}}\,\|\beta_\Sigma\|_{L^\infty}^2\,\|\Gamma_\Sigma^{\mathbf b}\|_{L^2}^2 ,
\end{equation}
which is~\eqref{eq:DriftStability} with $\delta_{\mathbf b}=C\varepsilon_{\mathbf b}/\sqrt{\underline\lambda}$ provided the second term is absorbed by the $-\tfrac14\|\beta_\Sigma\|_{L^\infty}^2\|\Gamma_\Sigma^{\mathbf b}\|_{L^2}^2$ of~\eqref{eq:DriftStability}, which holds for $C\varepsilon_{\mathbf b}/\underline\lambda^{3/2}<1/4$, i.e.\ for $\varepsilon_{\mathbf b}$ sufficiently small compared to $\underline\lambda^{3/2}$.

\begin{remark}[Honesty about the perturbative regime]\label{rem:PerturbativeHonesty}
The bound~\eqref{eq:MeanCurvatureSmallness} is the load-bearing input of Step 2: it converts the MOTS condition $H_\Sigma+\tr_\Sigma k=0$ into the smallness of $\tr_\Sigma k$ in $L^\infty$. This is the point where the perturbative theorem uses the full product-cylinder closeness, not merely the smallness of $\mathrm{tf}_\Sigma\,k$ that enters~\eqref{eq:CollarSmallness} through $\beta_\Sigma$. Crucially, $\beta_\Sigma$ (built from $\mathrm{tf}_\Sigma\,k$) and $\tr_\Sigma k$ are independent components of $k|_\Sigma$; the MOTS condition is what relates them, and the perturbative closeness is what makes the relation uniform. Outside this perturbative regime, the small-$|\beta_\Sigma|$ region is the locus where Hypothesis~\ref{hyp:DriftStability} carries genuine nontrivial content, and where the present argument does not close.
\end{remark}

\textbf{Step 3 (commutator absorption).}
The commutator bound~\eqref{eq:CommutatorBound} gives, by the same Young inequality,
\begin{equation}\label{eq:CommutatorAbsorb}
    |C_\Sigma|\;\le\;\tfrac{C\,\varepsilon_{\mathbf b}}{\sqrt{\underline\lambda}}\,\mathfrak q_{L_\Sigma}[\Gamma_\Sigma^{\mathbf b}]
    \;+\;\tfrac14\,\|\beta_\Sigma\|_{L^\infty}^2\,\|\Gamma_\Sigma^{\mathbf b}\|_{L^2}^2 .
\end{equation}
Combining~\eqref{eq:FirstForm},~\eqref{eq:TransportVsCoercivity}, and~\eqref{eq:CommutatorAbsorb} yields
\[
    \int_\Sigma(\tr_\Sigma k)\,\Gamma_\Sigma^{\mathbf b}\,dA
    \;\ge\;
    \bigl(1-\tfrac{C\,\varepsilon_{\mathbf b}}{\sqrt{\underline\lambda}}\bigr)\,\mathfrak q_{L_\Sigma}[\Gamma_\Sigma^{\mathbf b}]
    \;>\;0
\]
for $\varepsilon_{\mathbf b}<\tfrac12\sqrt{\underline\lambda}/C$, which is~\eqref{eq:PerturbativeDeltaB}.
\end{proof}

\subsection{Sign-free Penrose from drift-anchor stability}

\begin{theorem}[Drift-anchored Green-trace sign]\label{thm:DriftAnchoredGreenSign}
Assume the hypotheses of Theorem~\ref{thm:PerturbativeDriftStability}. Then the drift-anchored Green-trace pairing is nonnegative:
\begin{equation}\label{eq:DriftAnchoredGreenSign}
    \int_\Sigma(\tr_\Sigma k)\,\Gamma_\Sigma^{\mathbf b}\,dA\;\ge\;0 .
\end{equation}
\end{theorem}

\begin{proof}
Combine the decomposition of Lemma~\ref{lem:DriftAnchorDecomposition} with the bounds of Theorem~\ref{thm:PerturbativeDriftStability}. The drift transport and commutator are absorbed by $L_\Sigma$-coercivity and the negative-definite anchor potential $-\tfrac14\|\beta_\Sigma\|_{L^\infty}^2\|\Gamma_\Sigma^{\mathbf b}\|_{L^2}^2$ for $\varepsilon_{\mathbf b}$ sufficiently small, yielding~\eqref{eq:DriftAnchoredGreenSign}.
\end{proof}

\begin{theorem}[Spacetime Penrose inequality under drift-stability]\label{thm:PenroseDriftAnchor}
Let $(M,g,k)$ be an asymptotically flat initial data set satisfying the dominant energy condition. Let $\Sigma$ be a closed future-trapped surface for which the sealed Jang construction yields a strictly stable MOTS blow-up with $\lambda_1(L_\Sigma)>0$, and suppose the collar geometry satisfies the smallness~\eqref{eq:CollarSmallness} of Theorem~\ref{thm:PerturbativeDriftStability}. Then
\[
    \sqrt{\frac{A(\Sigma)}{16\pi}}\;\le\;m_{\mathrm{ADM}}(g) .
\]
\end{theorem}

\begin{proof}
By Theorem~\ref{thm:DriftAnchoredGreenSign} the drift-anchored Green-trace sign~\eqref{eq:DriftAnchoredGreenSign} holds. The interface contribution to the ADM mass coefficient is, by Lemma~\ref{lem:InterfaceADMGreen} applied to the drift-anchored potential (the lemma is invariant under the change $\Gamma_\Sigma\mapsto\Gamma_\Sigma^{\mathbf b}$ because the anchored correction is a solution of a strictly stable elliptic equation with bounded source, hence $o(1)$ in $H^1$ relative to $\Gamma_\Sigma$),
\[
    \delta m_\Sigma\;=\;\frac1{8\pi}\int_\Sigma(\tr_\Sigma k)\,\Gamma_\Sigma^{\mathbf b}\,dA\;+\;o(1)\;\ge\;o(1).
\]
The $o(1)$ remainder is the standard smoothing contribution of Section~\ref{app:SmoothingDetails}, which is nonnegative by the $L^{3/2}$ negative-part bound. Combining the interface sign with the Bray--Khuri identity of Section~\ref{app:BK_Identity} gives $\phi\le1$ and hence the spacetime Penrose inequality as in Theorem~\ref{thm:PenroseFromWeightedRiccati}.
\end{proof}

\subsection{Honest assessment and comparison with the prior Green-trace route}

We are explicit about the relationship of this section to the existing Green-trace calibration of Section~\ref{sec:GreenTraceCalibration} and to the no-go theorem of Section~\ref{app:LogicalStructure}.

\begin{remark}[What is genuinely new]\label{rem:DriftAnchorNew}
The structural innovation of this section is the identification of the drift one-form $\beta_\Sigma$ in~\eqref{eq:DriftOneForm} as a carrier of the trace data that the Green-trace sign requires, and the algebraic identity of Lemma~\ref{lem:DriftAnchorDecomposition} that converts the Green-trace pairing into a transport term, a sign-definite anchor potential, and a controlled commutator. The pointwise favorable-jump hypothesis of the direct MOTS argument required $\tr_\Sigma k\ge0$ everywhere; the Green-trace calibration of Section~\ref{sec:GreenTraceCalibration} required the conjectural weighted Riccati positivity of Conjecture~\ref{cor:PenroseFromDualCalibration}; the present route requires instead Hypothesis~\ref{hyp:DriftStability}, which is an integrated, drift-weighted inequality and is established in the perturbative regime by Theorem~\ref{thm:PerturbativeDriftStability}.
\end{remark}

\begin{remark}[What remains open]\label{rem:DriftAnchorOpen}
Hypothesis~\ref{hyp:DriftStability} is proven here only in the perturbative regime~\eqref{eq:CollarSmallness}, i.e.\ for collars that are $C^{2,\alpha}$-close to a product cylinder with $\mathrm{tf}_\Sigma\,k$ uniformly small in $W^{1,\infty}$. Outside this regime the drift-stability inequality~\eqref{eq:DriftStability} is a nontrivial constraint that relates the tangential drift $\beta_\Sigma$, the trace $\tr_\Sigma k$, and the anchored Green potential $\Gamma_\Sigma^{\mathbf b}$. We do not currently have a proof of~\eqref{eq:DriftStability} in the general non-perturbative regime; the non-time-symmetric input to the spacetime Penrose inequality along the Jang-equation route therefore remains conditional on either (i) Hypothesis~\ref{hyp:DriftStability} in full generality, (ii) the weighted Riccati positivity conjecture of Section~\ref{sec:GreenTraceCalibration}, or (iii) the pointwise favorable-jump hypothesis. The drift-anchor decomposition of Lemma~\ref{lem:DriftAnchorDecomposition} is, however, unconditional, and it identifies the precise analytic obstruction that any of (i)--(iii) must address.
\end{remark}

\begin{remark}[Comparison with the no-go theorem]\label{rem:VsNoGo}
Theorem~\ref{thm:NoGoVariational} rules out a \emph{bare} adjoint-inverse upgrade from the adjoint-cone Euler condition to the pointwise sign $\tr_\Sigma k\ge0$. The drift-anchor route does not attempt such an upgrade; it leaves the pointwise sign unset and instead routes the Green-trace pairing through the drift transport, where the MOTS stability operator enters with a coercive sign. The no-go theorem is therefore not in conflict with the present construction; it is precisely the obstruction that motivates the change of channel from the bare adjoint-inverse to the drift-anchored transport-diffusion.
\end{remark}
